\section{Fixed angular twists}
\label{sec:angular}

The first-moment comparison also holds after insertion of a fixed
angular Fourier mode. The primary-generator convention matters here:
it permits every integer mode, with no divisibility condition.
Throughout this appendix, $\ell\in\mathbb Z$ is fixed before $X$ tends
to infinity. All implied constants and subsequent fixed parameter
choices may depend on $\ell$.
For nonzero $z\in\mathbb C$, put
\[
 \theta_\ell(z)=(z/|z|)^\ell,\qquad
 G_\ell(n)=\theta_\ell(n)G(n),\qquad
 \mathcal D_\ell(n)=\theta_\ell(n)
       \bigl(G(n)-c_*\mu^2(n)N(n)^{-1/6}\bigr).
\]
As before, ideal sums use the unique primary generators.

\begin{theorem}[Comparison in each fixed angular mode]
\label{thm:angular}
For every fixed integer $\ell$,
\begin{equation}
 \sum_{\substack{N\pi\le X\\\pi\equiv1\pmod3\ \mathrm{prime}}}
 \theta_\ell(\pi)\bigl(G(\pi)-c_*N(\pi)^{-1/6}\bigr)
 =o_\ell\!\left(\frac{X^{5/6}}{\log X}\right).
 \label{eq:angular-comparison}
\end{equation}
In particular, for every fixed nonzero integer $\ell$,
\begin{equation}
 \sum_{\substack{N\pi\le X\\\pi\equiv1\pmod3\ \mathrm{prime}}}
       \theta_\ell(\pi)G(\pi)
 =o_\ell\!\left(\frac{X^{5/6}}{\log X}\right).
 \label{eq:angular-cancellation}
\end{equation}
\end{theorem}

The assertion is pointwise in the integer $\ell$. It makes no claim
uniform in a growing angular parameter. We prove the additional
analytic interfaces and then verify their use in the prime
decomposition and sharp-cutoff argument.

\subsection{Hecke characters with a fixed infinity type}

For an ideal $\mathfrak a$ prime to $3$, define
$\Theta_\ell(\mathfrak a)=\theta_\ell(a)$, where $a$ is its primary
generator. Products of primary generators are primary, so
$\Theta_\ell$ is multiplicative. If $\alpha$ is any generator prime
to $3$ and $u(\alpha)$ is the unit making $u(\alpha)\alpha$ primary,
then
\[
 \Theta_\ell((\alpha))
       =u(\alpha)^\ell(\alpha/|\alpha|)^\ell.
\]
The unit factor depends only on $\alpha\bmod3$. Thus $\Theta_\ell$
is a unitary Hecke character of angular order $\ell$, with finite
conductor dividing $(3)$. Its nonzero infinity type prevents it from
being principal when $\ell\ne0$. Multiplying by a finite-order cubic
character changes no angular order. If $\chi$ has finite conductor
$\mathfrak f_\chi$, the primitive character $\Psi$ inducing
$\chi\Theta_\ell$ has conductor dividing
$\operatorname{lcm}(\mathfrak f_\chi,(3))$; in particular,
$D=N\mathfrak f_\Psi\le9N\mathfrak f_\chi$.

\begin{lemma}[Fixed-infinity Hecke inputs]\label{lem:angular-hecke}
Put $k=|\ell|/2$. The complete primitive $L$-function has completion
\begin{equation}\label{eq:angular-hecke-fe}
 \Lambda(z,\Psi)=\left(\frac{\sqrt{3D}}{2\pi}\right)^z
       \Gamma(z+k)L(z,\Psi),\qquad
 \Lambda(z,\Psi)=\epsilon_\Psi\Lambda(1-z,\overline\Psi),
 \quad |\epsilon_\Psi|=1.
\end{equation}
If $\ell\ne0$, or if $\ell=0$ and $\chi$ is nonprincipal cubic,
the completed function is entire of order one. In these cases, for
every fixed $A,M>0$,
\begin{equation}\label{eq:angular-prime-sw}
 \sum_{N\mathfrak p\le x}\Psi(\mathfrak p)\log N\mathfrak p
       \ll_{\ell,A,M}x(\log x)^{-M},
 \qquad N\mathfrak f_\chi\le(\log x)^A.
\end{equation}
Ramified primes have character value zero. Removing additional
Euler factors is permitted with their explicit finite-prime error.
\end{lemma}

\begin{proof}
Rajan~\cite[Section 1, pp.\ 281--283]{Rajan1998} gives the
unitary completion, contragredient functional equation and
entire-order-one assertion. There is one complex place, field
discriminant $3$, angular parameter $\ell$ and no imaginary norm
parameter. Its gamma factor and conductor factor give exactly
\eqref{eq:angular-hecke-fe}. The dual angular order is $-\ell$.
A nonzero angular component cannot be removed by multiplying by a
finite character or an imaginary norm power: both are trivial on
the connected archimedean unit circle, whereas its $\ell$th angular
character is nontrivial for $\ell\ne0$. Hence $\Psi$ is not
principal or norm-equivalent to a real quadratic character when
$\ell\ne0$. For $\ell=0$, a nonprincipal cubic character is nonreal.
There is therefore no exceptional real zero in either case.

For clarity we give the uniformity deduction for
\eqref{eq:angular-prime-sw}. Apply the published
\cite[Theorem 1.2]{KanekoThorner2025}, which explicitly restates
the prime-number estimate of \cite[Theorem 5.13]{IwaniecKowalski2004}.
Viewed over rational primes, $L(z,\Psi)$ has degree $m=2$ and
local parameters of modulus at most one. In particular its
logarithmic-derivative coefficients satisfy $|a(p^j)|\le2$, and
the required second-moment hypothesis follows from
$\sum_{n\le x}|a(n)|^2\Lambda(n)^2\le4x(\log x)^2$.
Gamma duplication gives conductor $q=3D$, parameters
$\mu_1=k$, $\mu_2=k+1$, and analytic conductor
$\mathfrak C=3D(k+3)(k+4)$. Absolute convergence is supplied by
the Euler product, and the functional-equation hypothesis by
\eqref{eq:angular-hecke-fe}. The zero-free-region hypothesis is
Rajan's Proposition 1; the same conductor and infinity-parameter
dependence is recorded in
\cite[Lemma 2.1]{ChenParkSwaminathan2015}.
The absence of an exceptional character just proved permits the
bound
\[
 x\bigl(\log(x\mathfrak C)\bigr)^4
   \exp\!\left(-\frac{c\log x}
                         {\log\mathfrak C+\sqrt{\log x}}\right)
       +O(\sqrt x).
\]
This is $O_{\ell,A,M}(x\log^{-M}x)$ for every $M$ in the stated
conductor range. Removing higher prime-ideal powers from the
logarithmic derivative costs $O(\sqrt x\log^2x)$ and proves
\eqref{eq:angular-prime-sw}.
\end{proof}

Partial summation removes the logarithmic prime weight and permits
a norm twist $(N\mathfrak p)^{it}$ of logarithmic height, at cost
$O(1+|t|)$ absorbed in the freely chosen saving. The same argument
permits smooth norm weights and interval restrictions. We always
apply \eqref{eq:angular-hecke-fe} to the complete primitive ideal
sum. In particular, if the primitive character is unramified above
$3$, its omitted Euler factor is first restored; the value at that
prime is its actual Hecke value, not a chosen element phase.

\subsection{The angular Voronoi formula and Type I estimates}

We use the all-integer-index form of
\cite[Propositions 5.2--5.3]{DR2024}. For primary squarefree $r$,
put $R=Nr$ and define
\[
 \mathcal C_{r,\ell}(V;W)=
 \sum_{\substack{d,u\\(d,r)=1}}|d|G(ru)
       \theta_\ell(rud^3)W\left(\frac{N(ud^3)}V\right).
\]
In the notation and normalization of
Theorem~\ref{thm:voronoi-background}, the exact formula is
\begin{align}
 \mathcal C_{r,\ell}(V;W)
 ={}&\boldsymbol1_{\ell=0}A_0\mathcal MW(5/6)V^{5/6}
                  \frac{\varphi(r)}{R^{7/6}}\notag\\
 &+\frac{R^{1/2}\theta_{3\ell}(r)}{3^{7/2}(2\pi)^2}
 \sum_{\substack{\nu\ne0,\ d\\(d,r)=1}}
 \frac{a_r(\nu)b_r(\nu)\theta_{-\ell}(d^3\nu)}
      {N\nu\,(Nd)^{5/2}}\,
 \check W_\ell\left(\frac{(2\pi)^4N(d^3\nu)V}{R^2}\right),
 \label{eq:angular-voronoi}
\end{align}
where the support and absolute coefficient bounds are exactly those
in that theorem, and
\[
 \check W_\ell(v)=\frac1{2\pi i}\int_{(-\sigma)}
    v^z Q_\ell(z)\mathcal MW(z)\,dz,\qquad
 Q_\ell(z)=
 \frac{\Gamma(5/6+|\ell|/2-z)\Gamma(7/6+|\ell|/2-z)}
      {\Gamma(z+|\ell|/2-1/6)\Gamma(z+|\ell|/2+1/6)}.
\]
Here $0<\sigma<1/10000$ suffices. To check the phases, multiply
\cite[Equation (5.80)]{DR2024} by $G(r)\theta_\ell(r)$.
The outside angular factor becomes
$\theta_\ell(r)(\bar r/r)^{-\ell}=\theta_{3\ell}(r)$.
The coefficient formula (5.79), with the theta index
$\lambda^{-3}\nu$, supplies the displayed dual phase and the same
absolute bounds. The main-term indicator removes any angular phase
from the residue.

On $\Rea z=-\sigma$, Stirling's formula gives
$|Q_\ell(z)|\ll_{\ell,\sigma}(1+|\Im z|)^{2+4\sigma}$.
Consequently the absolute-convergence argument in
Theorem~\ref{thm:voronoi-background} proves
\begin{equation}
 \mathcal C_{r,\ell}(V;W)
 =\boldsymbol1_{\ell=0}A_0\mathcal MW(5/6)V^{5/6}
        \frac{\varphi(r)}{R^{7/6}}
       +O_{\ell,\epsilon}(X^\epsilon R^{1/2})
 \label{eq:angular-completed-low}
\end{equation}
under its same polynomial scale bounds, fixed compact support and
logarithmic derivative bounds.

The estimate at large heights requires more than the absence of an
angular residue. We first bound the completed sum: when its norm
length is short, an ordinary mean value suffices; otherwise the
Voronoi formula bounds its dual frequency range. This supplies the
angular replacement for the radial metaplectic mean-square input.

\begin{lemma}[Completed angular sums at large heights]
\label{lem:angular-completed-height}
Let $r$ be primary and squarefree, put $R=Nr$, and let $W$ have
fixed compact support in $(0,\infty)$ and uniformly bounded
derivatives. For fixed $\ell$, define
\[
 \mathcal C_{r,\ell,t}(V;W)=
 \sum_{\substack{d,u\\(d,r)=1}}|d|G(ru)
       \theta_\ell(rud^3)N(ud^3)^{it}
       W\left(\frac{N(ud^3)}V\right).
\]
Suppose $T\ge2$, $V\gg1$, and $R,V,T$ are bounded by fixed
powers of $X$. Write
\[
 M_{r,\ell,t}(V;W)=
 \boldsymbol1_{\ell=0}A_0\mathcal MW(5/6+it)V^{5/6+it}
                   \frac{\varphi(r)}{R^{7/6}}.
\]
Then
\begin{equation}
 \int_{|t|\asymp T}
       |\mathcal C_{r,\ell,t}(V;W)-M_{r,\ell,t}(V;W)|
           \frac{dt}T
 \ll_{\ell,\epsilon}X^\epsilon\sqrt V\,R^{1/4}\sqrt T.
 \label{eq:angular-completed-height}
\end{equation}
\end{lemma}

\begin{proof}
If $V\le\sqrt R\,T^2$, collect equal rational norms of $ud^3$.
Divisor multiplicities bound the squared coefficient sum by
\[
 X^\epsilon\sum_d Nd\,
       \#\{u:N(ud^3)\asymp V\}
       \ll X^\epsilon V\sum_d(Nd)^{-2}
       \ll X^\epsilon V.
\]
The ordinary mean-value inequality bounds the mean absolute size
of the full sum by $X^\epsilon\sqrt V\sqrt{1+V/T}$, which is at
most the right side of \eqref{eq:angular-completed-height}.
The possible main term is
$O_D(V^{5/6}R^{-1/6}T^{-D})$. Its ratio to the claimed bound is
at most $R^{-1/4}T^{1/6-D}\le1$ for $D\ge1$, by the current
inequality $V\le\sqrt R\,T^2$. Thus the pole-subtracted claim
also holds in this branch.

Suppose now $V>\sqrt R\,T^2$. In the angular transform use
$W(x)x^{it}$ and the outside scalar $V^{it}$; its Mellin factor is
$\mathcal MW(z+it)$. On $\Rea z=-A$, the gamma quotient grows
as $(1+|\Im z|)^{2+4A}$ for fixed $\ell,A$.
For any sufficiently small fixed $\delta>0$, truncate at
\[
N(d^3\nu)\le I_{\max}=X^\delta R^2T^4/V.
\]
This cutoff uses the dyadic height $T$ and is independent of $t$.
The omitted tail is smaller than any prescribed negative power
of $X$ by taking $A$ sufficiently large. Indeed the transform
contributes the power $(R^2T^4/(VN(d^3\nu)))^A$; the extra
$T^2$ and the absolutely summed coefficients are absorbed by
$X^{-\delta A}$ after all polynomial scale bounds are fixed.
Rapid decay in the translated Mellin variable controls its remaining
tail. Partition the retained finite sum into norm dyads $I$, so
\[
 I\ll X^\delta R^{3/2}T^2.
\]
More explicitly, put $D_0=R^2T^4/V$. The absolutely summed
coefficients with the additional weight $N(d^3\nu)^{-\sigma}$
cost $O_{\sigma,\eta}(R^\eta)$. For $A>\sigma$ the discarded
part is therefore
$O_{\ell,A,\sigma,\eta}(R^{1/2+\eta}T^2D_0^\sigma
X^{-\delta(A-\sigma)})$.
All factors preceding the negative power have fixed polynomial
bounds, justifying the stated choice of $A$.

Move each finite transform to $\Rea z=1/2$. Its first numerator
pole has real part $5/6+|\ell|/2>1/2$, and
\[
 |Q_\ell(1/2+iy)|=1.
\]
Translating the imaginary variable makes the weight
$\mathcal MW(1/2+iv)$ and the norm exponent $i(v-t)$.
The gamma quotient and conductor phase are scalars for the whole
frequency polynomial. Minkowski's inequality thus uses a common
rapidly decreasing $L^1$ majorant independent of $t$.
On this line the absolute coefficient, before summing frequencies,
is a fixed constant times
\[
 \sqrt V\,R^{-1/2}
       \frac{|a_r(\nu)b_r(\nu)|}{Nd\sqrt{N\nu}}.
\]

Write the dual support as
$\nu=\lambda^j\zeta h y(h')^3$ and fix $d,j,\zeta,h,h'$.
Put $D_d=Nd$, $E=Nh$, $P'=Nh'$, and
$g=N((\lambda\nu,r))\le\min(R,EP')$.
The free squarefree variable has norm length
$Y'\ll I/(D_d^3 3^j E(P')^3)$.
The coefficient bounds and the ordinary mean value, after extracting
$\sqrt V$, give
\[
 \ll_{\ell,\epsilon}X^\epsilon R^{-1/2}
 \frac{3^{-j/6}g}{D_d P'\sqrt E}
       \left(1+\sqrt{Y'/T}\right).
\]
Here the free coefficients have size $O((Ny)^{-1/2})$;
collecting equal norms costs only divisor factors. The endpoint
$j=-1$ changes constants.
For the first part use $g\le\sqrt{REP'}$. The $j$ sum is
geometric, the $d$ sum is logarithmic, the $h$ choices cost at
most $\tau(r)$, and
$\sum_{h'\mid r^\infty}(Nh')^{-1/2}\ll_\epsilon R^\epsilon$.
Its total is $O(X^\epsilon)$.
For the second part use $g\le EP'$; after substituting $Y'$ its
total is at most
\[
 X^\epsilon\sqrt{\frac I{RT}}
 \sum_{d,j,h,h'}\frac{3^{-2j/3}}
                       {D_d^{5/2}(P')^{3/2}}
 \ll X^\epsilon\sqrt{\frac I{RT}}
 \ll X^\epsilon R^{1/4}\sqrt T.
\]
The sum over $h$ is charged by a divisor bound; the other displayed
sums converge or cost $R^\epsilon$. Reselecting the small losses
proves \eqref{eq:angular-completed-height}.

\end{proof}

\begin{proposition}[Angular Type I estimates]
\label{prop:angular-typeI}
Suppose $RU\asymp X$, $1\le R\ll X^{.51}$, and $(a_r)$ is
divisor-bounded on $Nr\asymp R$. If $W_r$ has common fixed compact
support in $(0,\infty)$ and each fixed derivative is bounded by a
fixed logarithmic power uniformly in $r$, then
\begin{equation}
 \sum_r a_r\sum_u\mathcal D_\ell(ru)W_r(Nu/U)
       \ll_\ell X^{5/6-1/100}.
 \label{eq:angular-typeI-low}
\end{equation}
For the high-height assertion suppose instead that the derivative
bounds are fixed independently of $X$. For
$2\le T\ll X^{1/6+\rho}$ and every fixed $D>0$,
\begin{align}
 \int_{|t|\asymp T}\sum_r|a_r|
 \left|\sum_u G_\ell(ru)(Nu)^{it}W_r(Nu/U)\right|\frac{dt}T
 \ll_{\ell,\epsilon}{}&
       X^{1/2+\epsilon}R^{3/4}T^{1/2}\notag\\
 &+\boldsymbol1_{\ell=0}
       O_D(X^{5/6}L^C T^{-D}),
 \label{eq:angular-typeI-high}
\end{align}
where $\rho$ is the fixed parameter of \eqref{eq:power-parameters}
and $C$ depends only on the coefficient bounds.
\end{proposition}

\begin{proof}
Nonsquarefree levels vanish. For squarefree $r$ the exact inverse
completion is
\begin{equation}
 \sum_u G_\ell(ru)W_r(Nu/U)
 =\sum_{(c,r)=1}\mu(c)|c|\theta_{3\ell}(c)
       \mathcal C_{r,\ell}(U/(Nc)^3;W_r).
 \label{eq:angular-inversion}
\end{equation}
Indeed, after $e=cd$, the angular factor is
$\theta_\ell(rue^3)$ and the divisor coefficient is
$|e|\sum_{c\mid e}\mu(c)$. This is zero unless $e=1$.
In particular, the sign of the angular index on $c$ is positive.

For low heights take the cutoffs of Proposition~\ref{prop:typeI}:
$C_*=X^{.13}$ for $R\le X^{.40}$ and $C_*=X^{.02}$ otherwise.
The short errors from \eqref{eq:angular-completed-low} total
$O_{\ell,\epsilon}(X^\epsilon R^{3/2}C_*^{3/2})
 \ll X^{.795+\epsilon}$.
When $\ell=0$ the main term and its tail are exactly those already
computed there. When $\ell\ne0$ the Voronoi main term is absent.
The model is small in this case for a separate lattice reason.
For a fixed nonzero integer $\ell$,
\begin{equation}
 \sum_{(u,r)=1}\mu^2(u)\theta_\ell(u)W(Nu/U)
       \ll_{\ell,\epsilon}X^\epsilon\sqrt U.
 \label{eq:angular-lattice}
\end{equation}
To see this, a smooth annular function
$\theta_\ell(z)W(Nz/Y)$ has zero area integral and lattice
discrepancy $O_\ell(L^C(1+\sqrt Y))$, by comparison on fundamental
cells and integration of its gradient. Square-divisor and
coprimality inversion give $u=c^2dv$, with $(c,r)=1$ and $d\mid r$.
Rotation contributes only $\theta_\ell(c^2d)$. For
$Nc\ll\sqrt U$, the constant errors cost
$O(\sqrt U\,\tau(r))$; the radius errors cost at most
$O(\sqrt U\log U)$ times a divisor factor of $r$.
Terms whose scaled support contains no nonzero lattice point vanish.
This proves \eqref{eq:angular-lattice}. Only the area integral
vanishes exactly, not the discrete angular sum.

Apply this estimate to $x^{-1/6}W_r(x)$. After the outer sum, the
model error is
\[
 O_{\ell,\epsilon}(X^\epsilon R^{5/6}U^{1/3})
       \ll X^{.589+\epsilon}.
\]
For the long completion tail, combining $c,d$ into $e$ leaves
$\theta_{3\ell}(e)\theta_\ell(r)\theta_\ell(u)$ as independent
modulus-one coefficient factors. The same counting bound
$X^{1+\epsilon}E^{-3/2}$ and the same cubic-sieve bound
\[
 X^{1/2+\epsilon}
       \left(R+U/E^3+(X/E^3)^{2/3}\right)^{1/2}
\]
therefore apply on $Ne\asymp E$. The first gives exponent $.805$
in the small-$R$ branch; in the other branch the three exponents
are $.755$, $.77$, and $5/6-.02$. These, the short errors, and the
radial main-term tail when $\ell=0$ prove
\eqref{eq:angular-typeI-low}.

For high heights, let $b_0$ be a common upper endpoint of the
supports of $W_r$. A nonempty completed sum has
$Nc\le(b_0U)^{1/3}$ and hence $V=U/(Nc)^3\ge1/b_0$.
We may therefore apply Lemma~\ref{lem:angular-completed-height}
inside the inverse completion. With the norm twist its coefficient is
$\mu(c)|c|\theta_{3\ell}(c)(Nc)^{3it}$.
At $V=U/(Nc)^3$ its error is
$X^\epsilon\sqrt U\,(Nc)^{-1}R^{1/4}\sqrt T$.
Only $Nc\ll U^{1/3}$ can contribute, so this sum is logarithmic.
The main term is absent for $\ell\ne0$. For $\ell=0$ it is bounded
by $O_D(U^{5/6}R^{-1/6}T^{-D})$ after inversion: the Mellin
factor has arbitrary height decay and the $c$ sum has exponent $2$.
Finally, summing $|a_r|$ gives the error
$X^{1/2+\epsilon}R^{3/4}\sqrt T$ and the stated pole bound.
This proves \eqref{eq:angular-typeI-high}.
\end{proof}

\subsection{Smooth duality and the exceptional moments}

Call $\beta_\ell(b)=\Theta_\ell(b)\beta_0(b)$ an angular prime
convolution if $\beta_0$ is exactly an admissible convolution in
\eqref{eq:prime-convolution}. Its independent radial prime weights,
squarefree restriction, lower prime bound, bounded number of factors
and derivative hypotheses are all retained.

\begin{proposition}[Fixed-angular exceptional moments]
\label{prop:angular-dual}
Proposition~\ref{prop:dual} holds with $\beta$ replaced by
$\beta_\ell$. The constants $\delta,\eta$ and the implied constant
may additionally depend on the fixed integer $\ell$.
\end{proposition}

\begin{proof}
We first prove the version of Lemma~\ref{lem:smooth-character-duality}
needed for a full smooth factor. Let $\Psi$ be the primitive
nonprincipal character inducing $\chi\Theta_\ell$ in the
dominant-full-sum reduction, where $\chi$ is finite-order cubic,
including its fixed local factors. There is no additional imaginary
norm parameter: all norm oscillation is the explicit factor
$(N\mathfrak n)^{it}$. Mellin inversion gives
\[
 \sum_{\mathfrak n}\Psi(\mathfrak n)(N\mathfrak n)^{it}
       W(N\mathfrak n/Z)
 =\frac1{2\pi i}\int \mathcal MW(z)Z^zL(z-it,\Psi)\,dz.
\]
The sum is over all integral ideals. Move to $\Rea z=-A$ and
apply \eqref{eq:angular-hecke-fe}. For fixed $k=|\ell|/2$,
the quotient $\Gamma(1-z+it+k)/\Gamma(z-it+k)$ grows with
degree $1+2A$ in $1+|\Im z-t|$. Absolute convergence on this
line therefore allows truncation into dual norm dyads
\begin{equation}\label{eq:angular-smooth-cutoff}
 J\ll_{\ell,\epsilon}Y^\epsilon D(1+|t|)^2/Z,
\end{equation}
with an arbitrarily small power tail. Within each row average fix
$Z,t,W$ and the local and Euler-divisor case, as in the radial
Mellin reduction. Choose one cutoff using the largest permitted
conductor and the same smooth dyadic partition for every row.
This ensures that neither a row-dependent cutoff nor a varying
norm shift enters the coefficients.

Move each retained finite dyad to $\Rea z=1/2$. The first
numerator gamma pole has real part $1+k$, so none is crossed.
Writing $z=1/2+i\tau$ and $C_D=\sqrt{3D}/(2\pi)$, the
row-dependent factor is
\[
 \epsilon_\Psi C_D^{2i(t-\tau)}
 \frac{\Gamma(1/2+k-i(\tau-t))}
      {\Gamma(1/2+k+i(\tau-t))}.
\]
It has modulus one. The dyad amplitude is $\sqrt{Z/J}$, its
character is $\overline\Psi$, and
$|\mathcal MW(1/2+i\tau)|$ supplies a common rapidly decreasing
$L^1$ majorant independent of conductor and height. Its mass is
$O(Y^\epsilon)$ under the stipulated logarithmic derivative
bounds. Minkowski's inequality introduces only the common
additional norm shift $\tau$. On primary dual arguments the
angular multiplier is the common coefficient $\Theta_{-\ell}$.

The local factors require some care. Fix the small twists and
the finitely many conductor possibilities above $3$ as in the
dominant-full-sum argument of Section~\ref{sec:dual}. Remove
the missing Euler factors before applying the primitive equation.
When the primitive character is unramified at $\lambda$, split
its dual $\lambda$ powers and use $\overline\Psi(\lambda)^j$.
For each fixed $j$ this is a row scalar of modulus at most one;
on the remaining argument the angular coefficient is common.
It need not equal the raw element phase
$\theta_{-\ell}(\lambda)^j$. The same finite exclusions and
subpower divisor costs as in the radial proof suffice. For
$\ell\ne0$ nonprincipality follows from the infinity type; for
$\ell=0$ it follows from the surviving large-prime part of the
row conductor used there.

We now check the entire moment argument. Repeated-prime and
prime-power discrepancies have unchanged support and absolute
divisor bounds. Thus both discrepancy estimates of
Section~\ref{sec:dual} apply. If $T_\ell f(n)=\Theta_\ell(n)f(n)$,
multiplicativity gives the exact identity
$T_\ell(f*g)=(T_\ell f)*(T_\ell g)$. In particular the short
convolution identity used there becomes
\[
 T_\ell\Lambda=
 \bigl(2T_\ell m-(T_\ell m)*(T_\ell m)*(T_\ell\mathbf1)\bigr)
       *T_\ell(\log N).
\]
All truncation ranges and norm dyads are unchanged. In a product
of factors, including repeated copies, the merged coefficient is
multiplied by $\Theta_\ell$ of the product argument. The angular
order is not multiplied by the number of factors. The resulting
coefficients are still common to all cubic-character rows and
divisor-bounded.

If a length exponent tends to one, the corresponding full smooth
factor is treated by \eqref{eq:angular-smooth-cutoff}. The bound
$|t|\le Y^{.36}$, with the internal allowance $Y^{.37}$, yields
the same dual length $J\le Y^{.74+o(1)}$ and second-moment
exponents $2.16$ in the first configuration and
$2/3+2(1/3+.37)$ in the second. Both are strictly below $7/3$.
Thus this case is removed exactly with a power gap. Every remaining
sieve is applied to the finite cubic row character with
$\Theta_\ell$ in its arbitrary common coefficients.

With full smooth factors excluded, every limiting length exponent is
strictly less than one. The remaining moment argument in the proof
of Proposition~\ref{prop:dual} applies to the twisted polynomials:
their coefficients are common to all rows and divisor-bounded, and
these properties survive the fixed products and repetitions used
there. The large-value inequality \eqref{eq:dual-large-values}
therefore holds unchanged for a set of $Y^{r+o(1)}$ rows.
In the second configuration the same
ordinary-sieve argument, repeating one positive factor, gives
$r\le\max(4/3,z)-2z/3<2/3$, with $1<z<2$, and the required
contradiction.

In the first configuration the multiplicity inequality
\eqref{eq:dual-multiplicity}, in the notation of that proof, is still
\[
 r\le\frac23-2\sum_i k_i(v_i-5a_i/6),
 \qquad \frac12\le\sum_i k_i a_i\le2.
\]
The radial multiplicity comparison thus produces a polynomial of
length $Z=Y^{z+o(1)}$, with $4/3<z<2$, whose values are
$Y^{5z/6+o(1)}$ on $Y^{2/3-o(1)}$ distinct rows.
The Gram matrix of Lemma~\ref{lem:character-gram} is unchanged:
$\Theta_\ell$ belongs to the common coefficient vector, and cancels
against its conjugate in every row product. The radial weight stays
radial. Writing $R$ for the number of these rows, after the same
division by $ZR$ the left exponent is
$2z/3$, whereas the diagonal and off-diagonal exponents are
\[
 z-\frac23,\qquad
 \frac{z+\max(1,2-2z/3)-2/3}{2}.
\]
Their gaps are $(2-z)/3$ and, respectively, $(3z-4)/6$ for
$z\le3/2$ or $(z-1)/6$ for $z\ge3/2$. Each is positive.
Choosing the small-power losses after the fixed multiplicity gives
the same contradiction. Keep $\ell$ fixed in the uniform-neighborhood
argument at the end of that proof. A sequence of shrinking
neighborhoods and savings would again have a convergent subsequence
of length and value exponents; the selected multiplicity and losses
are fixed before taking the limit. The contradiction therefore gives
the required fixed $\delta,\eta>0$.
\end{proof}

\subsection{Twisted sequence conditions and corrected dispersion}

The condition \textup{(SW)} in Section~\ref{sec:dispersion} must
be checked for the twisted sequence itself. It is not an abstract
consequence of \textup{(SW)} for an arbitrary untwisted sequence.
For the actual decomposition coefficients we have the following.

\begin{lemma}[The angular sequence condition]
\label{lem:angular-sw}
Every nonsparse ordinary coefficient $\beta_0$ constructed in
Proposition~\ref{prop:decomposition} satisfies \textup{(SW)}
after replacement by $\beta_\ell=\Theta_\ell\beta_0$.
\end{lemma}

\begin{proof}
Repeat the largest-prime selection of
Lemma~\ref{lem:decomposition-sw}. Discard the part whose prime
factors all have norm at most $\exp(\sqrt{\log X})$ by its
absolute Rankin bound. For the rest, fix all factors except a
largest prime $\pi$, of dyadic norm length
$P\ge\exp(\sqrt{\log X})$, starting this dyadic subdivision
at that threshold. Their angular factor is a scalar.
The remaining prime sum is tested against
$\chi(\pi)\Theta_\ell(\pi)$, with $\chi$ nonprincipal cubic
of logarithmic conductor, a logarithmic norm height, and the
same smooth weights and norm intervals from the stopping rule.
Since $(\log X)^C\le(\log P)^{2C}$, the conductor is within
Lemma~\ref{lem:angular-hecke}. If $\ell\ne0$ the infinity type
prevents principality; if $\ell=0$ the nonprincipal cubic
condition does so. Partial summation supplies arbitrary logarithmic
saving uniformly in the remaining norm intervals and heights.
The polynomial-size exclusion integer deletes only $O(\log X)$
possible primes. This costs a logarithmic power per fixed
complement, harmless since $P\ge\exp(\sqrt{\log X})$.
Choose the prime saving after the fixed coefficient multiplicities,
norm-height powers and desired saving. Summing the complementary
factors then gives precisely the bound in \textup{(SW)}.
\end{proof}

\begin{proposition}[Angular dispersion and bilinear estimates]
\label{prop:angular-dispersion}
The estimates of Proposition~\ref{prop:dispersion} hold for
$\beta_\ell$, assuming \textup{(SW)} for $\beta_\ell$ wherever
it was required. For $B>X^{.40}$, the stipulated structured
class is exactly the angular prime-convolution class above.
The conclusions of Propositions~\ref{prop:bilinear-low}
and~\ref{prop:bilinear-height} consequently hold with the kernel
multiplied by $\theta_\ell(ab)$ and with all their other
hypotheses retained. The absolute small-$B$ parameter $\gamma$
is unchanged; other constants may depend on fixed $\ell$.
\end{proposition}

\begin{proof}
Keep $W_A(a)$ radial and put
\[
 F_{\ell,u}(a)=\sum_b\beta_\ell(b)(Nb)^{iu}G(b)
                       \overline{\chi_b(a)},\qquad
 S_\ell(u)=\sum_b\beta_\ell(b)(Nb)^{iu-1/6}.
\]
In every opened pair, the extra factor is
$\Theta_\ell(b_1)\overline{\Theta_\ell(b_2)}$ in the
coefficients. At $b_i=fln_i$ the common factor cancels, leaving
the separate twists on $n_1,n_2$. Poisson summation and its
plane transform still act on a radial weight. The diagonal,
large-gcd removal, transform tails and all general-coefficient
sieves and height means therefore have their original bounds.

At the two exceptional large-conductor configurations the length
argument still forces $f=l=1$, $B=X^{1/2+o(1)}$, small extra
twists and exclusions, and the internal height at most $B^{.36}$.
Proposition~\ref{prop:angular-dual} supplies their saving.
At small nonprincipal conductor use \textup{(SW)} for
$\beta_\ell$ itself. When $B\le X^{.40}$ the exceptional
configurations remain excluded by the original absolute length
gaps. Thus the power-width derivative bounds, sharp norm-cell
restrictions and unrestricted translated height in that range
retain the same $\gamma$, independently of $\ell$.

The cube frequencies give the main term
\begin{equation}\label{eq:angular-cube-main}
 c_*^2|S_\ell(u)|^2
       \sum_a\mu^2(a)W_A(a)(Na)^{-1/3}.
\end{equation}
This expression need not vanish for $\ell\ne0$. Its cancellation
uses the corrected square, as in the radial argument. Indeed the
mixed term before multiplication by $\overline{S_\ell(u)}$ is
\[
 \sum_{a,b}\beta_\ell(b)(Nb)^{iu}
                  W_A(a)(Na)^{-1/6}G(ab).
\]
Here the angular factor is only in the arbitrary outer
$b$-coefficient. The inner $a$-weight is radial, so the original
Proposition~\ref{prop:typeI} gives
\[
 c_*\sum_{a,b}\beta_\ell(b)(Nb)^{iu-1/6}
          \mu^2(ab)W_A(a)(Na)^{-1/3}
   +O\bigl(A^{-1/6}X^{5/6-\eta}\bigr).
\]
After multiplication by
$|S_\ell(u)|\ll B^{5/6}L^{O(1)}$, the error is
$O(\mathcal QX^{-\eta}L^{O(1)})$.
Roughness and the unchanged absolute coefficient moments remove
cross-coprimality at the original logarithmic cost. The positive
square, mixed term and model square thus have the same main
term \eqref{eq:angular-cube-main} with coefficients $1,-2,1$.
This proves the corrected dispersion bound.

In the bilinear sum use
$\alpha_\ell(a)=\Theta_\ell(a)\alpha_0(a)$ and
$\beta_\ell(b)$. The first phase disappears under Cauchy,
and its coefficient norm is unchanged. Restoring the model
gives exactly $\theta_\ell(ab)c_*\mu^2(ab)N(ab)^{-1/6}$.
All cross-coprimality errors use absolute values. In particular,
the three-prime transition still has saving $L^{-3/2}$.

For the integrated estimate, at large $B$ the diagonal remains
$A\sqrt B\ll X^{.80}$ and the other term has arbitrary
logarithmic saving. At small $B$ partition both norm variables
into the same $O(J)$ cells. Their coefficient norms $a_i,b_j$
are unchanged. Near pairs give
$X^{5/6}L^{O(1)}H^{-1/6}J^{-1/2}$ as in
\eqref{eq:near-cell-diagonal}; all absolute errors sum to
$X^{5/6}L^{O(1)}J^{3/2}\varepsilon_J^{1/2}$ as in
\eqref{eq:all-cell-error}. For far pairs, integration by parts
differentiates the height cutoff. The angular phases stay in
the coefficients, and the common Mellin shifts remain independent
of the averaging height. Thus the factor $(J/T)^q$ and the
unrestricted-shift use of dispersion are preserved. With
$J=L^{D'}$ or $X^\gamma$ the same positive gaps prove the
integrated estimate.
\end{proof}

\subsection{The exact decomposition and sharp cutoff}

We finally verify that these interfaces prove
Theorem~\ref{thm:angular}. The identities of
Section~\ref{sec:decomposition} apply to both squarefree-supported
kernels $G_\ell$ and $\mathcal D_\ell$. The prime detector,
distinguished tuples and semiprime subtraction are scalar
identities. Every stopping predicate depends on norms alone.
For a fixed stopping bin and its counts, multiplicativity
allocates the phase to independent coefficients
$\Theta_\ell(a)\alpha(a)$ and $\Theta_\ell(b)\beta(b)$;
the binomial weight is unchanged. The permitted overlaps are
still annihilated by the separate squarefree restriction on
each kernel. There is no new coupled cutoff or height dependence
in the coefficients.

No-stop terms have an angular inner variable and use
Proposition~\ref{prop:angular-typeI}. Every sparse support count,
absolute coefficient moment, Rankin bound and general-coefficient
cubic-sieve application is unchanged. For $B>X^{.40}$ the
coefficient is precisely $\Theta_\ell\beta_0$ with the original
independently smooth prime weights; Proposition~\ref{prop:angular-dual}
therefore applies. Lemma~\ref{lem:angular-sw} proves the required
sequence condition. The number of pieces, transition indices
and sparse remainders is exactly the one already counted.

Fix $\ell$ first. Choose the absolute small-$B$ $\gamma$,
then $\kappa,\rho$ as in \eqref{eq:power-parameters}, then
$\xi$ and the fixed tuple complexity. Subsequent neighborhood
sizes, logarithmic savings and constants may depend on $\ell$.
They are chosen in the order of Section~\ref{sec:perron}.
There is no dependence returning from an angular exceptional
moment to the absolute $\gamma$.

\begin{proof}[Proof of Theorem~\ref{thm:angular}]
Apply Lemma~\ref{lem:perron-truncation} to the prime coefficients
with their angular factors and $T_{\max}=X^{1/6+\rho}$.
Equal-norm collection has the same absolute divisor bound,
so the truncation error is
$O_\epsilon(X^\epsilon(1+X/T_{\max}))$
and is $o(X^{5/6}/\log X)$ for $\epsilon<\rho$.
At low heights the exact decomposition just verified, the
angular Type I comparison and the angular bilinear comparison
give the same arbitrary logarithmic savings. The three-prime
transition contributes
$O_\ell(X^{5/6}L^{-3/2}(1+\log L)^C)$, which is also little-oh
at the required scale. Sparse terms retain their power savings.

At high heights first bound the undecomposed angular prime model.
If $d_{\ell,n}$ is its coefficient after collecting equal norms,
then $\sum_n|d_{\ell,n}|^2\ll X^{2/3}L^{O(1)}$ by absolute
divisor multiplicity. The ordinary mean value therefore gives
the same $T^{-1/2}$ saving used in Section~\ref{sec:perron}.
This step does not use angular prime cancellation.
For the Gauss-sum term the integrated bilinear estimate is
Proposition~\ref{prop:angular-dispersion}. The high-height
Type I estimate gives the same two endpoint powers,
$.805+\epsilon$ and $5/6-3\kappa/8+\rho/2+\epsilon$.
Choose $\epsilon$ below the two fixed gaps to $5/6$.
Its pole term is absent unless $\ell=0$, when its arbitrary
height decay is the original one. The two-cell and far-pair
savings proved above pay for the same polynomial number of
pieces and the $O(L)$ height dyads. Hence the dyadic prime
comparison \eqref{eq:prime-dyadic-comparison} holds with
$\theta_\ell(\pi)$ inserted. Geometric dyadic summation gives
\eqref{eq:angular-comparison}.

For $\ell\ne0$, take $\chi$ principal in
Lemma~\ref{lem:angular-hecke} and remove the prime logarithm by
partial summation. Thus
$P_\ell(y)=\sum_{N\pi\le y}\theta_\ell(\pi)
 \ll_{\ell,M}y(\log y)^{-M}$ for every fixed $M$.
Abel summation now gives
\[
 \sum_{N\pi\le X}\theta_\ell(\pi)(N\pi)^{-1/6}
 =X^{-1/6}P_\ell(X)
       +\frac16\int_2^X P_\ell(y)y^{-7/6}\,dy
 =o_\ell(X^{5/6}/\log X).
\]
The lower endpoint contributes only a bounded term. Combining
this with \eqref{eq:angular-comparison} proves
\eqref{eq:angular-cancellation}.
\end{proof}

For $\ell=0$ the model is the radial one and retains coefficient
$(6/5)c_*$. Inert primes still have total size $O(\sqrt X)$.
The rational half-sum consequence \eqref{eq:rational-main}
is unchanged. For $\ell\ne0$ a conjugate pair instead contributes
$2\Rea(\theta_\ell(\pi)G(\pi))$; no identification with the
untwisted rational Kummer summand is asserted.

\subsection{Fixed powers of the Gauss-sum values}

The angular result also determines the fixed power moments whose
exponents are not divisible by three. In the normalization
\eqref{eq:gauss-definition}, the cubic Gauss-sum identity
\cite[Equation (1.8)]{DR2024} gives
\[
 G(\pi)^3=-\frac{\pi}{|\pi|}=-\theta_1(\pi),
 \qquad |G(\pi)|=1.
\]
Consequently, for every fixed integer $k$ with $3\nmid k$,
\begin{equation}\label{eq:gauss-power-cancellation}
 \sum_{\substack{N\pi\le X\\\pi\equiv1\pmod3\ \mathrm{prime}}}
       G(\pi)^k
 =\boldsymbol1_{k\in\{-1,1\}}\frac65c_*
       \frac{X^{5/6}}{\log X}
   +o_k\!\left(\frac{X^{5/6}}{\log X}\right).
\end{equation}
Indeed, for $k=3m+1$ and $k=3m-1$, respectively,
\[
 G(\pi)^{3m+1}=(-1)^m\theta_m(\pi)G(\pi),
 \qquad
 G(\pi)^{3m-1}=(-1)^m
                  \overline{\theta_{-m}(\pi)G(\pi)}.
\]
For $m\ne0$, Theorem~\ref{thm:angular} gives cancellation; for
$m=0$, Theorem~\ref{thm:main} and its complex conjugate give the
same real main term. Negative powers are well-defined because
$|G(\pi)|=1$.

Thus \eqref{eq:gauss-power-cancellation} proves the fixed-power
cases $3\nmid k$, $|k|>1$, of the complementary conjecture in
\cite[Equation (1.9)]{DR2024}. For a nonzero multiple $k=3m$,
the same identity instead gives
$G(\pi)^k=(-1)^m\theta_m(\pi)$. The Hecke prime estimates used
here do not bound these prime-angle sums at the stronger scale
$X^{5/6}/\log X$, so those cases remain outside the conclusion.
